\renewcommand{\headline}{Penta$\cdot$ludus - Latine}
\renewcommand{\tocent}{Latine}
\renewcommand{\tocline}{Regulae Latinae}
\renewcommand{\translator}{Dr.~J\"org Grimm, JS}
\renewcommand{\general}{ 
    \lettrine[lines=3]{L}{udus}
    mirabilis atque explicabilis velox. 
    Duo lusores ludent tertiam partem horae, tres aut quattuor ludent horam. 
    
    Alea non opus est. 
    Aptus lusoribus a quinque annis.
    Tabula pentagrammiforma, quattuor cohortes quinorum peditum coloratorum, quina impedimenta nigra et cana Penta\-ludum componunt.  Fit Ioh.~\noun{Aridus.}
}
\renewcommand{\choosext}{Cohortes}
\renewcommand{\choosex}{
    Lusor quisque agit pedites quinos, quibus est eadem forma, sed sunt diversi coloris. 
    
    Sic tua cohors peditem caeruleum, peditem rubrum et sequentes continet.
}

\renewcommand{\setupt}{Dispositio} 
\renewcommand{\setup}{ 
    In quinque nodis positis in circulo, qui circumvenit pentagramma, ponite pedites eiusdem coloris: pedites albos in nodo albo, caeruleos
    in nodo caeruleo, et sequentes. 
    
    Praeter illa decem aedicula quae nodos formant, existunt alia octoginta aedicula in viis, ubi pedites in ludo etiam possunt collocari. 
    
    Sunt etiam impedimenta nigra et im\-pe\-di\-men\-ta cana. 
    Impedimenta nigra ponite in nodis centri. In centro retinete impedimenta cana. 
}

\renewcommand{\objectivet}{Propositum} 
\renewcommand{\objective}{ 
    Destinatio per unum quemque peditem, qui occupat nodum externi circuli, est nodus eiusdem coloris in interno pentagrammatis. 
    Albo est eundum ad album, rubro ad rubrum, et sequentes. 
    
    Lusor victor qui duxit tres pedites ad nodum aequi coloris.
}

\renewcommand{\rulest}{Regulae } 
\renewcommand{\rules}{ 
    Duc peditum tuorum unum in quamvis directionem sequens aut viam circuli aut vias stellae, quoad per regulas potes occupare aediculum. 

    \myskip
    
    In quolibet aediculo libero, ubi  duae lineae conveniunt, tibi licet de via recta declinare sine retentione. 
    
    \myskip
    
    At numquam tibi licet transcendere nec impedimenta nec pedites alios.
    
    \myskip
    
    Potes vero ingredi in aediculum occupatum:
    
    \myskip
    
    Si ibi stat impedimentum nigrum, id pone in aediculo libero ad libitum. 
    
    Si ibi iam stat pedes alius, pedites positiones suas mutare debent.
    
    Si ibi sunt plures pedites (non potest fieri, nisi initio ludi), elige unum mutandum.
    
    \myskip
    
    Ita potes etiam mutare positiones duorum peditum tuorum.
    
    \myskip
    
    Eundem motum non licet iterum.
    
    \myskip
    
    Pedes progressus ad finem ludo abit. Colloca in centrum.
    
    Accipe inde impedimentum canum et colloca in aediculo libero, quoad tibi paret. 
    
    \myskip
    
    Si captas tale impedimentum canum, recolloca in centrum.
    
    \myskip
    
    Si non satis impedimenta cana, in ludo recolloca tale.
    
    \myskip
    
    Si alius lusor ad finem posivit peditem tuum, tibi is deinde removendus est.
    
    \myskip
    
    Orbis ultimus perficiatur.
}
